\problem{3.11}{Momentum distribution of the oscillator ground state}
Set $\ell=\sqrt{\hbar/(m\omega)}$. The normalized
position wavefunction is
$\psi_0(x)=(\pi\ell^2)^{-1/4}
\exp[-x^2/(2\ell^2)]$.
Fourier transformation with the $p$-normalization gives
\begin{align}
\Phi_0(p,t)&=\frac{\ee^{-\ii\omega t/2}}
{\sqrt{2\pi\hbar}}\int_{\mathbb R}
(\pi\ell^2)^{-1/4}
\ee^{-x^2/(2\ell^2)-\ii px/\hbar}\dd x\\
&=\boxed{(\pi m\hbar\omega)^{-1/4}
\ee^{-p^2/(2m\hbar\omega)}
\ee^{-\ii\omega t/2}}.
\end{align}
The classical momentum limit at $E_0=\hbar\omega/2$
is $p_c=\sqrt{2mE_0}=\sqrt{m\hbar\omega}$.
Therefore, with $u=p/\sqrt{m\hbar\omega}$,
\begin{equation}
P(|p|>p_c)=\frac2{\sqrt\pi}\int_1^\infty
\ee^{-u^2}\dd u=\operatorname{erfc}(1)
\simeq\boxed{0.16\quad\text{(two significant figures)}}.
\end{equation}
\qnote{The classically forbidden momentum tail
is the momentum-space counterpart of the
oscillator's position-space forbidden region.}

\problem{3.12}{Fourier coefficients of a free wave packet}
For the free-particle expansion
$\Psi(x,t)=
(2\pi)^{-1/2}\int\phi(k)
\ee^{\ii kx-\ii\hbar k^2t/(2m)}\dd k$,
the symmetric momentum Fourier transform yields
\begin{equation}
\Phi(p,t)=\frac1{\sqrt{2\pi\hbar}}
\int\ee^{-\ii px/\hbar}\Psi(x,t)\dd x.
\end{equation}
Insert the expansion and use
$\int_{\mathbb R}\ee^{\ii(k-p/\hbar)x}\dd x
=2\pi\delta(k-p/\hbar)$:
\begin{equation}
\boxed{\Phi(p,t)=\frac1{\sqrt\hbar}
\phi(p/\hbar)\ee^{-\ii p^2t/(2m\hbar)}}.
\end{equation}
Taking the squared modulus removes the phase:
\begin{equation}
\boxed{|\Phi(p,t)|^2=
\hbar^{-1}|\phi(p/\hbar)|^2},
\end{equation}
independent of time. This also follows
from the free Hamiltonian commuting with momentum.

\problem{3.13}{The position operator in momentum space}
Starting from the inverse transform,
$\Psi(x)=(2\pi\hbar)^{-1/2}\int
\ee^{\ii px/\hbar}\Phi(p)\dd p$,
differentiate the plane-wave factor:
$x\ee^{\ii px/\hbar}=
-\ii\hbar\,\partial_p\ee^{\ii px/\hbar}$.
Integration by parts (vanishing boundary term) then gives
\begin{equation}
x\Psi(x)=\frac1{\sqrt{2\pi\hbar}}
\int\ee^{\ii px/\hbar}
\left(\ii\hbar\,\partial_p\Phi(p)\right)\dd p.
\end{equation}
Since Fourier transformation is unitary,
\begin{equation}
\boxed{\avg x=\int_{\mathbb R}
\Phi^*(p)\,\ii\hbar\frac{\partial\Phi}{\partial p}\dd p}.
\end{equation}
Thus $\hat x=\ii\hbar\partial_p$
in the momentum representation, whereas
$\hat p$ acts by multiplication by $p$.

\problem{3.14}{Commutator algebra and functions of position}
Expand $\comm{A+B}C=(A+B)C-C(A+B)$ to obtain
$\boxed{\comm{A+B}C=\comm AC+\comm BC}$.
Similarly,
\begin{align}
\comm{AB}C
&=ABC-CAB\\
&=A(BC-CB)+(AC-CA)B\\
&=\boxed{A\comm BC+\comm AC\,B}.
\end{align}
The canonical commutator gives
$\comm{x^n}p=
x^{n-1}\comm xp+
\comm{x^{n-1}}p\,x$.
By induction,
\begin{equation}
\boxed{\comm{x^n}p=\ii n\hbar\,x^{n-1}}.
\end{equation}
If $f(x)=\sum_{n=0}^\infty a_nx^n$
and formal termwise differentiation is valid
on a common dense domain,
\begin{equation}
\boxed{\comm{f(x)}p=
\sum_{n=1}^\infty a_n\ii n\hbar x^{n-1}
=\ii\hbar f'(x)}.
\end{equation}
Direct verification:
$f(-\ii\hbar\psi')+
\ii\hbar(f\psi)'=\ii\hbar f'\psi$.

\problem{3.15}{Position–energy uncertainty relation}
Apply Robertson's inequality to $A=x$ and $B=H$:
$\sigma_x\sigma_H\ge
\frac12|\avg{\comm{x}{H}}|$.
For $H=p^2/(2m)+V(x)$,
$\comm{x}{V(x)}=0$ and
\begin{equation}
\comm{x}{p^2}=\comm xp\,p+
p\comm xp=2\ii\hbar p.
\end{equation}
Consequently
\begin{equation}
\boxed{\sigma_x\sigma_H\ge
\frac{\hbar}{2m}|\avg p|}.
\end{equation}
In a stationary energy eigenstate,
$\sigma_H=0$, but also
$\dd\avg x/\dd t=\avg p/m=0$
(for a normalizable state with no boundary
current). Therefore $\avg p=0$
and the inequality becomes $0\ge0$.

\problem{3.16}{Noncommutativity forbids a complete common eigenbasis}
Assume two linear operators $A,B$ have
a complete basis of simultaneous eigenvectors
$u_n$, with $Au_n=a_nu_n$ and $Bu_n=b_nu_n$.
For every basis vector,
\begin{equation}
[A,B]u_n=(a_nb_n-b_na_n)u_n=0.
\end{equation}
By linearity, $[A,B]$ annihilates every
finite linear combination of basis vectors.
On a finite-dimensional space this gives
$[A,B]=0$ on the whole space. For
unbounded operators in an infinite-dimensional
Hilbert space, the same conclusion holds
on a common invariant dense span, subject
to the required domains. The contrapositive is
that a genuinely nonzero commutator excludes
a complete common eigenbasis.

\problem{3.17}{All normalized minimum-uncertainty position states}
Equality in Cauchy--Schwarz requires
$(p-p_0)\Psi=\ii c(x-x_0)\Psi$
for an appropriate real nonzero constant $c$.
In the position representation,
\begin{equation}
-\ii\hbar\Psi'(x)-p_0\Psi(x)
=\ii c(x-x_0)\Psi(x).
\end{equation}
Divide by $\Psi$ and integrate:
\begin{equation}
\frac{\Psi'}{\Psi}=\frac{\ii p_0}{\hbar}
-\frac{c}{\hbar}(x-x_0),
\quad
\Psi(x)=A\exp\left[
-\frac{c(x-x_0)^2}{2\hbar}+
\frac{\ii p_0x}{\hbar}\right].
\end{equation}
Normalizability requires $c>0$. Writing
$c=\hbar/(2\sigma_x^2)$ and normalizing yields
\begin{equation}
\boxed{\Psi(x)=
(2\pi\sigma_x^2)^{-1/4}
\exp\left[-\frac{(x-x_0)^2}{4\sigma_x^2}
+\frac{\ii p_0x}{\hbar}\right]}.
\end{equation}
Its distribution has
$\avg x=x_0$, $\sigma_x$ as specified,
$\avg p=p_0$, $\sigma_p=\hbar/(2\sigma_x)$,
hence $\sigma_x\sigma_p=\hbar/2$.
\qnote{A physically irrelevant global phase
may multiply the solution.}

\problem{3.18}{Time derivatives of four expectation values}
For time-independent $H$ and any operator $Q$,
\begin{equation}
\frac{\dd}{\dd t}\avg Q=
\frac{\ii}{\hbar}\avg{\comm H Q}
+\left\langle\frac{\partial Q}{\partial t}\right\rangle.
\end{equation}
(a) $Q=I$: $\comm HI=0$ gives
$\boxed{\dd\avg I/\dd t=0}$, norm conservation.
(b) $Q=H$: $\comm HH=0$ gives
$\boxed{\dd\avg H/\dd t=0}$, energy conservation
when $\partial H/\partial t=0$.
(c) $Q=x$: $\comm Hx=-\ii\hbar p/m$ gives
\begin{equation}
\boxed{\dd\avg x/\dd t=\avg p/m}.
\end{equation}
(d) $Q=p$: $\comm Hp=\ii\hbar V'(x)$ gives
\begin{equation}
\boxed{\dd\avg p/\dd t=-\avg{V'(x)}}.
\end{equation}
These are the Ehrenfest equations.

\problem{3.19}{Free and harmonic-oscillator motion of the mean}
For a free particle, $V=0$, hence
$\dd\avg p/\dd t=0$. Integrating twice gives
\begin{equation}
\boxed{\avg x(t)=\avg x(0)+
\frac{\avg p(0)}m t}.
\end{equation}
For $V(x)=m\omega^2x^2/2$, Ehrenfest's
equations imply
$m\dd^2\avg x/\dd t^2=-m\omega^2\avg x$.
With $x_0=\avg x(0)$ and $p_0=\avg p(0)$,
\begin{equation}
\boxed{\avg x(t)=x_0\cos\omega t+
\frac{p_0}{m\omega}\sin\omega t}.
\end{equation}
Both results hold for any normalized state
with the appropriate finite expectations.

\problem{3.20}{Energy–time inequality for a two-level box superposition}
Use the normalized state from Problem 2.5,
$(\phi_1\ee^{-\ii E_1t/\hbar}
+\phi_2\ee^{-\ii E_2t/\hbar})/\sqrt2$,
$E_n=n^2\hbar\Omega$,
$\Omega=\pi^2\hbar/(2ma^2)$.
Its energy uncertainty is
\begin{equation}
\sigma_H=\frac{|E_2-E_1|}{2}
=\boxed{\frac{3\hbar\Omega}{2}}.
\end{equation}
By integrating sine products,
$\avg x=a/2-\frac{16a}{9\pi^2}\cos(3\Omega t)$,
$\avg{x^2}=a^2(\frac13-\frac{5}{16\pi^2})
-\frac{16a^2}{9\pi^2}\cos(3\Omega t)$, so
\begin{equation}
\boxed{\sigma_x^2(t)=a^2
\left[\frac1{12}-\frac{5}{16\pi^2}
-\frac{256}{81\pi^4}\cos^2(3\Omega t)\right]}.
\end{equation}
Furthermore
\begin{equation}
\left|\frac{\dd\avg x}{\dd t}\right|=
\frac{16a\Omega}{3\pi^2}
|\sin(3\Omega t)|.
\end{equation}
With $\Delta t=\sigma_x/|\dd\avg x/\dd t|$
(where the denominator is nonzero),
\begin{equation}
\sigma_H\Delta t=
\frac{9\pi^2\hbar}{32a}
\frac{\sigma_x(t)}{|\sin(3\Omega t)|}
\ge\frac{\hbar}{2}.
\end{equation}
To check the last inequality algebraically,
square it and use
$\sigma_x^2/a^2\ge
(256/(81\pi^4))\sin^2(3\Omega t)$;
the difference is
$1/12-5/(16\pi^2)-256/(81\pi^4)>0$.
At instants when the mean position is momentarily
stationary, $\Delta t$ is infinite or undefined.

\problem{3.21}{Energy–time inequality for a traveling Gaussian}
For the freely moving Gaussian of Problem 2.42,
$\Psi(x,0)\propto\ee^{-ax^2+\ii \ell x}$,
where $a>0$, the momentum distribution is
Gaussian with mean $p_0=\hbar\ell$ and
variance $s_p^2=a\hbar^2$.
For a normal random variable,
$\mathrm{Var}(p^2)=4p_0^2s_p^2+2s_p^4$;
hence, with $H=p^2/(2m)$,
\begin{equation}
\boxed{\sigma_H=
\frac{\hbar^2}{m}\sqrt{a\ell^2+\frac{a^2}{2}}}.
\end{equation}
Writing $s=2\hbar at/m$, the wavepacket width
and mean velocity are
\begin{equation}
\boxed{\sigma_x(t)=\frac{\sqrt{1+s^2}}{2\sqrt a}},
\quad
\boxed{\frac{\dd\avg x}{\dd t}=\frac{\hbar\ell}{m}}.
\end{equation}
For $\ell\ne0$,
\begin{equation}
\sigma_H\Delta t=
\frac{\hbar}{2}\sqrt{1+\frac{a}{2\ell^2}}
\sqrt{1+s^2}\ge\frac{\hbar}{2}.
\end{equation}
If $\ell=0$, the mean does not move and this
particular characteristic-time definition
is not informative.

\problem{3.22}{Energy–time and position–energy uncertainty coincide}
For $Q=x$ the Mandelstam–Tamm timescale is
$\Delta t=\sigma_x/|\dd\avg x/\dd t|$.
Ehrenfest gives $\dd\avg x/\dd t=\avg p/m$,
so the energy–time inequality
$\sigma_H\Delta t\ge\hbar/2$ becomes
\begin{equation}
\frac{\sigma_H\sigma_x m}{|\avg p|}
\ge\frac\hbar2
\quad\Longleftrightarrow\quad
\boxed{\sigma_x\sigma_H\ge
\frac{\hbar|\avg p|}{2m}},
\end{equation}
which is exactly Problem 3.15.
When $\avg p=0$ the intermediate division
is not permitted; the direct commutator form
remains valid.
